\documentclass[11pt]{amsart}

\usepackage[T1]{fontenc}
\usepackage{amsmath,amssymb,amsthm,mathtools}
\usepackage{booktabs}
\usepackage{enumitem}
\usepackage[hidelinks]{hyperref}
\usepackage[margin=1.1in]{geometry}

\newtheorem{theorem}{Theorem}[section]
\newtheorem{proposition}[theorem]{Proposition}
\newtheorem{corollary}[theorem]{Corollary}
\newtheorem{lemma}[theorem]{Lemma}
\newtheorem{conjecture}[theorem]{Conjecture}
\theoremstyle{definition}
\newtheorem{definition}[theorem]{Definition}
\theoremstyle{remark}
\newtheorem{remark}[theorem]{Remark}

\newcommand{\Z}{\mathbb Z}
\newcommand{\cP}{\mathcal P}
\newcommand{\Gr}{\operatorname{Gr}}
\newcommand{\Hilb}{\operatorname{Hilb}}
\newcommand{\GL}{G_{\mathrm L}}
\newcommand{\GR}{G_{\mathrm R}}
\newcommand{\ol}[1]{\overline{#1}}

\title{The Bigrassmannian Minimality Conjecture:\\
A Source-Grounded Closure by Partition Conjugation}
\author{Co-Math research synthesis}
\date{July 16, 2026}

\begin{document}

\begin{abstract}
Reiner, Woo, and Yong conjectured that two explicit Schur-function generating
families for every type-$A$ bigrassmannian basic ideal are both minimal.  A
2024 theorem of St.~Dizier and Yong proves minimality for the lower-adjoined
family and remarks that the other case should follow by transposing the
argument.  We give the missing explicit transfer.  Partition conjugation
induces an integral graded isomorphism between complementary Grassmannian
presentations, carries the right-adjoined family to the published
lower-adjoined family for a conjugate bigrassmannian parameter tuple, and
preserves ideal generation and irredundancy.  This proves the original
two-family conjecture, including the $a=0$ boundary.  We also record the
accepted integral basis of indecomposables, its Gaussian-binomial Hilbert
series, auxiliary rank formulas, bounded computation checks, and the precise
Co-Math evidence and review boundary.  No universal transition matrix between
the two families is asserted or needed.
\end{abstract}

\maketitle
\tableofcontents

\section{Historical statement and status}

Let $X=\operatorname{Fl}_n(\mathbb C)$ be the complete flag variety.  For a
type-$A_{n-1}$ bigrassmannian permutation $v$, Reiner--Woo--Yong associate a
basic ideal $J_v\subset H^*(X;\Z)$.  Their Theorem~5.4 gives two finite
Schur-function generating families.  They state the following conjecture
\cite{RWY}.

\begin{conjecture}[Reiner--Woo--Yong minimality conjecture]
For every type-$A$ bigrassmannian $v=v_{r,s,t,n}$, both generating families
from Theorem~5.4 are minimal: no proper subfamily generates $J_v$.
\end{conjecture}

The original paper verified the conjecture computationally whenever
$r\leq4$ and $n-r\leq5$.  It also proposed the stronger structural
expectation that the module of indecomposables should be free of rank equal to
the common number of generators.  For the family $n=4m$ and $a=b=m$, the
conjecture forces at least
\[
  \binom{2m}{m}
  \sim \frac{4^m}{\sqrt{\pi m}}
  =\frac{\sqrt{2}^{\,n+2}}{\sqrt{\pi n}}
\]
generators.  Thus the conjecture supplies an exponential obstruction to
uniformly short presentations.

St.~Dizier--Yong prove minimality of the lower-adjoined family over $\Z$
\cite{SDY}.  Their Remark following the proof identifies the other family as
the transposed family and says that the arguments could be transposed.  The
purpose of the proof below is to make the simpler functorial conjugation
transfer completely explicit, including all parameter and boundary checks.

\section{Notation and the two families}

A nonidentity bigrassmannian permutation is indexed by integers satisfying
\begin{equation}\label{eq:bigrass-hyp}
  1\leq t\leq r,s\leq n,
  \qquad t>r+s-n.
\end{equation}
Put
\begin{equation}\label{eq:parameters}
  q=n-r,
  \qquad i=s-t+1,
  \qquad j=r-t+1,
\end{equation}
and
\begin{equation}\label{eq:ab}
  a=\min(q-i,r-j),
  \qquad b=\min(i,j).
\end{equation}
For a partition $\lambda$, write $\lambda'$ for its conjugate.  The rectangle
$q^r$ has $r$ rows of length $q$, and $i^j$ has $j$ rows of length $i$.

Let $\Lambda$ be the ring of symmetric functions over $\Z$, with Schur basis
$\{s_\lambda\}$.  Define
\[
  I_{r,q}
  =\bigoplus_{\lambda\nsubseteq q^r}\Z s_\lambda,
  \qquad
  A_{r,q}=\Lambda/I_{r,q}
  \cong H^*(\Gr(r,n);\Z).
\]
An overline denotes projection to $A_{r,q}$.  In the Grassmannian
presentation, the contracted basic ideal has the integral Schur basis
\begin{equation}\label{eq:J-basis}
  J_v
  =\bigoplus_{i^j\subseteq\lambda\subseteq q^r}
    \Z\ol{s}_\lambda.
\end{equation}
The two generating families are
\begin{align}
  \GL(v)
  &=\left\{
    \ol{s}_\mu:
    i^j\subseteq\mu\subseteq(i^j,b^a)
  \right\},\label{eq:GL}\\
  \GR(v)
  &=\left\{
    \ol{s}_\mu:
    i^j\subseteq\mu
    \subseteq((i+a)^b,i^{j-b})
  \right\}.
  \label{eq:GR}
\end{align}
Both are indexed by partitions in a $b\times a$ rectangle and therefore have
cardinality $\binom{a+b}{a}$.  We call them lower-adjoined and
right-adjoined, respectively.

The projection $X\to\Gr(r,n)$ induces a split inclusion on cohomology.  By
\cite[lines 207--211 in the arXiv source]{SDY}, a subfamily minimally
generates the Grassmannian contraction if and only if it minimally generates
the original flag-variety basic ideal.  It is consequently enough to work in
$A_{r,q}$.

\section{The conjugate parameter tuple}

\begin{lemma}[Admissibility and parameter exchange]\label{lem:parameters}
Define
\begin{equation}\label{eq:vee-tuple}
  (r^\vee,s^\vee,t^\vee,n)
  =(q,q-i+j,q-i+1,n)
  =(n-r,n-s,n-r-s+t,n).
\end{equation}
Then $v^\vee=v_{r^\vee,s^\vee,t^\vee,n}$ is bigrassmannian and its associated
parameters satisfy
\[
  (i^\vee,j^\vee,a^\vee,b^\vee)=(j,i,a,b).
\]
\end{lemma}

\begin{proof}
First $r,s<n$.  Indeed, if $r=n$, then
$t>r+s-n=s$, contradicting $t\leq s$; the argument for $s<n$ is symmetric.
Thus $q=n-r$ and $n-s$ are positive.

The strict integral inequality in \eqref{eq:bigrass-hyp} gives
$t^\vee=n-r-s+t\geq1$.  Moreover,
\[
  t^\vee\leq n-r=r^\vee \quad\Longleftrightarrow\quad t\leq s,
\]
and
\[
  t^\vee\leq n-s=s^\vee \quad\Longleftrightarrow\quad t\leq r.
\]
Finally,
\[
 t^\vee-(r^\vee+s^\vee-n)
 =(n-r-s+t)-(n-r-s)=t>0.
\]
Hence the transformed tuple satisfies every bigrassmannian inequality.

Direct substitution gives
\begin{align*}
 i^\vee&=s^\vee-t^\vee+1=r-t+1=j,\\
 j^\vee&=r^\vee-t^\vee+1=s-t+1=i,\\
 a^\vee
 &=\min(n-r^\vee-i^\vee,r^\vee-j^\vee)\\
 &=\min(r-j,q-i)=a,\\
 b^\vee&=\min(i^\vee,j^\vee)=b.
\end{align*}
\end{proof}

The same inequalities also give
\begin{equation}\label{eq:box-ineq}
  i\leq q,
  \quad j\leq r,
  \quad i+a\leq q,
  \quad j+a\leq r,
  \quad b\leq i,j.
\end{equation}
Thus all endpoint partitions in \eqref{eq:GL} and \eqref{eq:GR} lie in the
ambient rectangle, as do their transformed analogues.

\section{Integral partition conjugation}

Let $\omega:\Lambda\to\Lambda$ be the standard involution
\[
  \omega(s_\lambda)=s_{\lambda'}.
\]
It is a graded ring involution and permutes the integral Schur basis.

\begin{lemma}[Rectangle-ideal descent]\label{lem:omega-quotient}
For all $r,q\geq1$,
\[
  \omega(I_{r,q})=I_{q,r}.
\]
Consequently $\omega$ induces a graded integral ring isomorphism
\begin{equation}\label{eq:omega-quotient}
  \omega_{r,q}:A_{r,q}\xrightarrow{\sim}A_{q,r},
  \qquad
  \ol{s}_\lambda\longmapsto\ol{s}_{\lambda'}.
\end{equation}
\end{lemma}

\begin{proof}
A partition $\lambda$ lies in $q^r$ precisely when
$\lambda_1\leq q$ and $\ell(\lambda)\leq r$.  These conditions are equivalent
to $(\lambda')_1\leq r$ and $\ell(\lambda')\leq q$, namely
$\lambda'\subseteq r^q$.  Conjugation therefore bijects the Schur basis
outside $q^r$ with the Schur basis outside $r^q$.  The equality of ideals and
the claimed quotient isomorphism follow basiswise over $\Z$.
\end{proof}

\begin{lemma}[Basic ideals and endpoint intervals]\label{lem:omega-families}
Under \eqref{eq:omega-quotient},
\[
  \omega_{r,q}(J_v)=J_{v^\vee},
  \qquad
  \omega_{r,q}(\GR(v))=\GL(v^\vee).
\]
\end{lemma}

\begin{proof}
By \eqref{eq:J-basis}, conjugation sends the interval
$i^j\subseteq\lambda\subseteq q^r$ to
$j^i\subseteq\lambda'\subseteq r^q$.  Lemma~\ref{lem:parameters} identifies
the latter interval with the Schur basis of $J_{v^\vee}$.

For the generating families, the endpoint identities are
\begin{equation}\label{eq:endpoint-conjugation}
  (i^j)'=j^i,
  \qquad
  ((i+a)^b,i^{j-b})'=(j^i,b^a).
\end{equation}
Partition conjugation preserves containment, so it bijects the full interval
defining $\GR(v)$ with the interval defining $\GL(v^\vee)$.

If $a=0$, both upper endpoints equal their cores:
\[
  (i^j,b^0)=((i+0)^b,i^{j-b})=i^j,
\]
and similarly after exchanging $i$ and $j$.  Thus the argument includes the
singleton boundary case without an exception.
\end{proof}

\section{Closure of the original conjecture}

\begin{theorem}[Two-family minimality]\label{thm:minimality}
For every type-$A$ bigrassmannian $v=v_{r,s,t,n}$, both $\GL(v)$ and $\GR(v)$
minimally generate the basic ideal $J_v$, integrally.  This holds both in the
Grassmannian contraction and in the original flag-variety cohomology ring.
\end{theorem}

\begin{proof}
St.~Dizier--Yong prove that $\GL(w)$ minimally generates $J_w$ for every
bigrassmannian $w$ \cite[Theorem~1.2 and its concluding corollary]{SDY}.
Apply that theorem to $w=v$ to obtain minimality of $\GL(v)$.

Apply it again to the admissible bigrassmannian $v^\vee$ from
Lemma~\ref{lem:parameters}.  By Lemmas~\ref{lem:omega-quotient} and
\ref{lem:omega-families}, the inverse image of $J_{v^\vee}$ is $J_v$ and the
inverse image of $\GL(v^\vee)$ is $\GR(v)$.  A ring isomorphism preserves
generation.  If a proper subfamily of $\GR(v)$ generated $J_v$, its image
would be a proper generating subfamily of $\GL(v^\vee)$, contradicting the
published minimality theorem.  Hence $\GR(v)$ is minimal.

The split-inclusion equivalence cited at the end of Section~2 transfers both
minimality statements from the Grassmannian contraction to the original
flag-variety basic ideal.
\end{proof}

\begin{remark}
The proof uses no formula expressing one family linearly in terms of the
other inside a fixed Grassmannian quotient.  It instead moves to the
complementary Grassmannian quotient by a canonical integral ring isomorphism.
Thus a universal $\GL$-to-$\GR$ transition identity is neither assumed nor
proved.
\end{remark}

\section{The stronger indecomposables theorem}

Let
\[
  A_+=\bigoplus_{d>0}(A_{r,q})_d,
  \qquad Q_v=J_v/A_+J_v.
\]
The accepted Co-Math kernel certificate combines the published generation and
no-syzygy theorem with a full-coordinate integral descent.  It gives the
following strengthening of Theorem~\ref{thm:minimality}.

\begin{theorem}[Integral indecomposables basis]\label{thm:Q-basis}
There is a graded integral isomorphism
\begin{equation}\label{eq:Q-basis}
  Q_v\cong
  \bigoplus_{\rho\subseteq b^a}\Z e_\rho,
  \qquad
  \deg e_\rho=ij+|\rho|.
\end{equation}
The classes of $\GL(v)$ form one such homogeneous basis, and the classes of
$\GR(v)$ form another.  Consequently
\begin{equation}\label{eq:Hilbert}
  \Hilb_{Q_v}(z)
  =z^{ij}\binom{a+b}{a}_z,
\end{equation}
and $Q_v$ is torsion-free.
\end{theorem}

\begin{proof}[Proof recorded by the accepted certificate chain]
Index the lower family as
\[
  g_\rho=\ol{s}_{(i^j,\rho)},
  \qquad \rho\subseteq b^a.
\]
Published generation gives a graded integral surjection
\[
  \bigoplus_{\rho\subseteq b^a}\Z e_\rho\twoheadrightarrow Q_v,
  \qquad e_\rho\mapsto[g_\rho].
\]
A nonzero homogeneous kernel element in excess degree $N>0$ would give a
relation
\[
 \sum_{|\beta|=N}c_\beta g_\beta
 =\sum_{|\lambda|<N}f_\lambda g_\lambda,
 \qquad f_\lambda\in(A_{r,q})_{N-|\lambda|}.
\]
The left side is nonzero because the shapes $(i^j,\beta)$ are distinct Schur
basis elements in $q^r$.  This is exactly the homogeneous syzygy excluded by
the no-syzygy theorem of \cite{SDY}.  The accepted full-coordinate integral
descent certificate verifies that no coefficient is lost in passing to the
source linear system.  Thus the kernel vanishes for $N>0$.

For $N=0$, only $g_\varnothing=\ol{s}_{i^j}$ occurs.  Since $J_v$ is
generated in degrees at least $ij$, $(A_+J_v)_{ij}=0$, so its class is
nonzero.  If $a=0$, this is the entire argument because the indexing box has
only the empty partition.  This proves the lower basis statement.

The isomorphism of Lemma~\ref{lem:omega-quotient} maps positive-degree ideals
to positive-degree ideals and hence induces
\[
  J_v/A_+J_v\xrightarrow{\sim}
  J_{v^\vee}/A_+J_{v^\vee}.
\]
Lemma~\ref{lem:omega-families} transfers the lower basis for $v^\vee$ to the
right basis for $v$.  Formula \eqref{eq:Hilbert} follows by summing
$z^{ij+|\rho|}$ over partitions in $b^a$.
\end{proof}

Complementation in the box,
\[
  \rho^\vee=(b-\rho_a,\ldots,b-\rho_1),
\]
is an involution and $|\rho^\vee|=ab-|\rho|$.  Therefore
\[
  \langle e_\rho,e_\sigma\rangle
  =\delta_{\sigma,\rho^\vee}
\]
defines a perfect integral pairing
\[
  (Q_v)_d\times(Q_v)_{2ij+ab-d}\longrightarrow\Z.
\]
For $a=0$, the empty partition is declared self-complementary.

\section{Rank formulas accepted during the investigation}

Let
\[
  R_{a,b}(N)=\#\{\rho\subseteq b^a:|\rho|=N\}.
\]
Then Theorem~\ref{thm:Q-basis} gives
\[
  \operatorname{rank}(Q_v)_{ij+N}=R_{a,b}(N).
\]
If $p_a(m)$ is the number of partitions of $m$ with at most $a$ parts, with
$p_a(m)=0$ for $m<0$, an accepted inclusion--exclusion certificate gives
\begin{equation}\label{eq:rank-ie}
  R_{a,b}(N)
  =\sum_{S\subseteq\{1,\ldots,a\}}
    (-1)^{|S|}
    p_a\!\left(N-|S|b-\sum_{s\in S}s\right).
\end{equation}
Box complementation gives the symmetry
\[
  R_{a,b}(N)=R_{a,b}(ab-N).
\]

For the three-row box $b^3$, define $P(t)=0$ for $t<0$.  If
$t=6u+r\geq0$ with $0\leq r<6$, then
\[
P(t)=
\begin{cases}
3u^2+3u+1,&r=0,\\
3u^2+4u+1,&r=1,\\
3u^2+5u+2,&r=2,\\
3u^2+6u+3,&r=3,\\
3u^2+7u+4,&r=4,\\
3u^2+8u+5,&r=5.
\end{cases}
\]
Here $P(t)$ counts partitions of $t$ with at most three parts.  The accepted
piecewise quasi-polynomial certificate is
\[
R_{3,b}(N)=
\begin{cases}
0,&N<0\text{ or }N>3b,\\
P(N),&0\leq N\leq b,\\
P(N)-P(N-b-1)-P(N-b-2)-P(N-b-3),
  &b<N\leq\lfloor3b/2\rfloor,\\
R_{3,b}(3b-N),&\lfloor3b/2\rfloor<N\leq3b.
\end{cases}
\]

\section{Bounded computations and their proper role}

The proof of Theorem~\ref{thm:minimality} is symbolic and does not depend on
finite experimentation.  Co-Math nevertheless retained two useful audits.

\begin{enumerate}[label=\arabic*.]
\item The parameter-conjugation checker examined every bigrassmannian tuple
for $2\leq n\leq12$: $1001$ tuples total, including $495$ with $a>0$ and
$506$ with $a=0$.  Every asserted inequality, shape containment, conjugation
identity, and boundary convention passed.  The executable computation
artifact is
\texttt{dd8c3f05d598063b3da443dec00d07e26292748e8b66fd64e4cf2300d33856c9}.

\item A separate Smith-normal-form experiment used the interval
\[
 [2,2,2,0,0,0]\subseteq\lambda\subseteq[5,5,5,5,5,5]
\]
and all complete $h_1,\ldots,h_5$ Pieri rows in degrees $6$ through $30$.
Verified unimodular witnesses gave quotient ranks
\[
 (1,1,2,2,2,1,1)
\]
in degrees $6$ through $12$, zero thereafter through degree $30$, and no
torsion.  The computation artifact is
\texttt{6e0ed66127f15e5b06fbfd7c2c69c95905c664df4aaf1c92166348e3496fc5a5}.
\end{enumerate}

These computations are checks and examples only.  Neither was promoted to an
all-parameter theorem without a symbolic proof.

\section{Co-Math provenance and independent review}

The durable Co-Math state is preserved in the branch
\texttt{comath/research-exploration-mode} of
\url{https://github.com/HanzhangYin/pi-mono}, commit
\texttt{ace829d7}.  Content-addressed artifacts make the source excerpts,
claims, reports, criticism, and computations immutable.

\begin{center}
\small
\begin{tabular}{@{}lll@{}}
\toprule
Task & Role in the proof & Accepted report or ledger hash\\
\midrule
228 & Symbolic admissibility plus bounded audit
  & \texttt{251e640f38b88d9a5...}\\
229 & Endpoint shapes, boxes, and $a=0$
  & \texttt{3046a5a5e812cc7b...}\\
286 & Right-family indecomposables transfer
  & \texttt{b644ba099b1ae76f...}\\
296 & Lower-family graded-kernel certificate
  & \texttt{4a11210ef593a85f...}\\
301 & Integral quotient conjugation
  & \texttt{b9354e944eac34d9...}\\
326 & Three-row rank quasi-polynomial
  & \texttt{42c40bf8cea24fa5...}\\
328 & General inclusion--exclusion rank formula
  & \texttt{dfcb882173218746...}\\
334 & Consolidated theorem boundary
  & \texttt{c131f5897c63b586...}\\
336 & Source-grounding and missing-review diagnosis
  & \texttt{334dc58de7359d81...}\\
337 & Fresh end-to-end adversarial closure audit
  & \texttt{8f3dfb9330c77bdc...}\\
\bottomrule
\end{tabular}
\end{center}

For task~337, the specialist artifact is
\texttt{b722c59f0baf5dd996cc330e4eba9b9a3a477c14adb0cd7fab46a9b1b9f23eee};
the independent critic is
\texttt{d6ba1aac1f88f60379f0b472f9e169a8d54471c98a23302e1407c8bb8e5d30e6};
the final report is
\texttt{81e32be5ba7bd5dd5cce3f7c73f4cb245c9e4e331398dedef0e231837960a4bd};
and the independent skeptic verdict is
\texttt{d7f57938698ec321136dbf52345b00128dd7a07bf7481097aa8c78a991ae3cbb}.
The critic independently recomputed the admissibility inequalities, endpoint
conjugation, integral ideal descent, quotient-basis transfer, and
irredundancy.  The skeptic accepted the theorem with no concerns.

This review standard is analogous to a carefully source-checked mathematical
argument with independent model criticism.  It is not a proof-assistant
formalization: the definitions and deductions have not been encoded in Lean,
Coq, Isabelle, or another kernel-checked formal system.

\section{Claims deliberately left outside the theorem}

The following are not needed for, and are not claimed by, the closure proof.
\begin{itemize}
\item No parameter-uniform transition matrix between $\GL(v)$ and $\GR(v)$
inside the same quotient has been proved.
\item The universal transition identity explored during the research remains
a separate theorem candidate or conjecture.
\item The rejected larger transition computation lacked a certificate that
its selected relation lattice was the complete relation lattice; it was not
used.
\item Minimal generating sets for arbitrary type-$A$ Schubert ideals $I_w$
remain a broader open problem.
\item The result is not presently kernel-checked by a proof assistant.
\end{itemize}

\section{Conclusion}

The original mathematical target is closed without the optional transition
identity.  The 2024 theorem proves minimality of the lower-adjoined family.
The explicit admissible parameter involution \eqref{eq:vee-tuple}, together
with the integral partition-conjugation isomorphism
\eqref{eq:omega-quotient}, sends the right-adjoined family to a published
lower-adjoined family and transfers minimality back.  The same map transfers
the stronger integral indecomposables basis.  All edge cases, especially
$a=0$, are included.

\begin{thebibliography}{9}

\bibitem{RWY}
V.~Reiner, A.~Woo, and A.~Yong,
\emph{Presenting the cohomology of a Schubert variety},
Trans. Amer. Math. Soc. \textbf{363} (2011), 521--543,
\href{https://arxiv.org/abs/0809.2981}{arXiv:0809.2981}.

\bibitem{SDY}
A.~St.~Dizier and A.~Yong,
\emph{Presenting the cohomology of a Schubert variety: Proof of the
minimality conjecture},
J. London Math. Soc. \textbf{109} (2024), no.~1, e12832,
\href{https://doi.org/10.1112/jlms.12832}{doi:10.1112/jlms.12832},
\href{https://arxiv.org/abs/2209.02011}{arXiv:2209.02011}.

\end{thebibliography}

\end{document}
